% GENERATED FILE. Edit canonical lessons, then run npm run generate:books.
% canonical-source-hash: fnv1a64:f761509b8f1f29c8
% canonical-lessons: BN-C325-upodesh, BN-C325-porinam, BN-C325-pap, BN-C325-shodhrano, BN-C325-badhyo-hooya

\chapter{A Story's Lesson}
\label{ch:bn-a-story-s-lesson}
\hlchaptermodality{325}

\begin{chapteropening}
\textbf{By the end of this chapter:} \emph{I can say what a story teaches about right and wrong.}

Complete the last lesson of chapter 325: I can say what a story teaches about right and wrong.
\end{chapteropening}

\section[\texorpdfstring{\bn{উপদেশ} (upôdesh)}{উপদেশ (upôdesh)}]{\bn{উপদেশ} (upôdesh) — a moral lesson; counsel}

\label{lesson:BN-C325-upodesh}



\begin{quote}
Before the new one: say the Bengali for to coax, then the Bengali for to get away with it.
\end{quote}

\subsection*{You'll want to know: \bn{উপদেশ}}
\textbf{\bn{উপদেশ}} — \emph{upôdesh} — \textquotedblleft{}a moral lesson; counsel\textquotedblright{}.

\begin{quote}
\bn{গল্পের} \bn{শেষে} \bn{দাদু} \bn{একটা} \bn{উপদেশ} \bn{দিলেন।}
\end{quote}

\begin{quote}
\emph{gôlper sheshe dādu ekṭā upôdesh dilen.}
\end{quote}

\begin{quote}
At the end of the story my grandfather gave us a lesson to remember.
\end{quote}

\begin{cousinweb}
\textbf{\bn{উপদেশ}} (\emph{upôdesh}) is Sanskrit \emph{upadeśa}, 'pointing out; instruction, advice' (Monier-Williams); it has nothing to do with \textbf{\bn{দেশ}}, a country. Everyday advice is \textbf{\bn{পরামর্শ}}.
\end{cousinweb}

\begin{grammarlens}[title={What you've built}]
One more word for telling what happened.
\end{grammarlens}

\subsection*{Your turn}
\begin{itemize}
\raggedright
  \item \emph{Say it:} \emph{upôdesh}
  \item \emph{Say it:} \emph{upôdesh}, once more
  \item \emph{Say it:} \emph{gôlper sheshe dādu ekṭā upôdesh dilen.} — At the end of the story my grandfather gave us a lesson to remember.
  \item \emph{Recall:} say the Bengali for to coax, then the Bengali for to get away with it, then say \emph{upôdesh} again
\end{itemize}

\begin{tcolorbox}[breakable,colback=teal!4,colframe=teal!35!black,title={Before you move on}]
What does \bn{উপদেশ} mean? (\textquotedblleft{}A moral lesson; counsel\textquotedblright{}.) Say it once more.
\end{tcolorbox}
\section[\texorpdfstring{\bn{পরিণাম} (pôrinām)}{পরিণাম (pôrinām)}]{\bn{পরিণাম} (pôrinām) — an outcome; a consequence}

\label{lesson:BN-C325-porinam}



\begin{quote}
Before the new one: say the Bengali for to get away with it, then the Bengali for a moral lesson.
\end{quote}

\subsection*{You'll want to know: \bn{পরিণাম}}
\textbf{\bn{পরিণাম}} — \emph{pôrinām} — \textquotedblleft{}an outcome; a consequence\textquotedblright{}.

\begin{quote}
\bn{লোভের} \bn{পরিণাম} \bn{ভালো} \bn{হয়} \bn{না।}
\end{quote}

\begin{quote}
\emph{lobher pôrinām bhālo hôy nā.}
\end{quote}

\begin{quote}
Greed never ends well.
\end{quote}

\begin{cousinweb}
\textbf{\bn{পরিণাম}} (\emph{pôrinām}) is Sanskrit \emph{pariṇāma}, 'change, development, ripening' (Monier-Williams): what a thing turns into in the end.
\end{cousinweb}

\begin{grammarlens}[title={What you've built}]
One more word for telling what happened.
\end{grammarlens}

\subsection*{Your turn}
\begin{itemize}
\raggedright
  \item \emph{Say it:} \emph{pôrinām}
  \item \emph{Say it:} \emph{pôrinām}, once more
  \item \emph{Say it:} \emph{lobher pôrinām bhālo hôy nā.} — Greed never ends well.
  \item \emph{Recall:} say the Bengali for to get away with it, then the Bengali for a moral lesson, then say \emph{pôrinām} again
\end{itemize}

\begin{tcolorbox}[breakable,colback=teal!4,colframe=teal!35!black,title={Before you move on}]
What does \bn{পরিণাম} mean? (\textquotedblleft{}An outcome; a consequence\textquotedblright{}.) Say it once more.
\end{tcolorbox}
\section[\texorpdfstring{\bn{পাপ} (pāp)}{পাপ (pāp)}]{\bn{পাপ} (pāp) — a sin}

\label{lesson:BN-C325-pap}



\begin{quote}
Before the new one: say the Bengali for a moral lesson, then the Bengali for an outcome.
\end{quote}

\subsection*{You'll want to know: \bn{পাপ}}
\textbf{\bn{পাপ}} — \emph{pāp} — \textquotedblleft{}a sin\textquotedblright{}.

\begin{quote}
\bn{গল্পে} \bn{বলে} \bn{যে} \bn{মিথ্যা} \bn{বলা} \bn{পাপ।}
\end{quote}

\begin{quote}
\emph{gôlpe bôle je mitthā bôlā pāp.}
\end{quote}

\begin{quote}
The story says that telling lies is a sin.
\end{quote}

\begin{cousinweb}
\textbf{\bn{পাপ}} (\emph{pāp}) is Sanskrit \emph{pāpa}, 'bad, evil', and as a noun 'a sin' (Monier-Williams).
\end{cousinweb}

\begin{grammarlens}[title={What you've built}]
One more word for telling what happened.
\end{grammarlens}

\subsection*{Your turn}
\begin{itemize}
\raggedright
  \item \emph{Say it:} \emph{pāp}
  \item \emph{Say it:} \emph{pāp}, once more
  \item \emph{Say it:} \emph{gôlpe bôle je mitthā bôlā pāp.} — The story says that telling lies is a sin.
  \item \emph{Recall:} say the Bengali for a moral lesson, then the Bengali for an outcome, then say \emph{pāp} again
\end{itemize}

\begin{tcolorbox}[breakable,colback=teal!4,colframe=teal!35!black,title={Before you move on}]
What does \bn{পাপ} mean? (\textquotedblleft{}A sin\textquotedblright{}.) Say it once more.
\end{tcolorbox}
\section[\texorpdfstring{\bn{শোধরানো} (shodhrāno)}{শোধরানো (shodhrāno)}]{\bn{শোধরানো} (shodhrāno) — to mend one's ways; to put right}

\label{lesson:BN-C325-shodhrano}



\begin{quote}
Before the new one: say the Bengali for an outcome, then the Bengali for a sin.
\end{quote}

\subsection*{You'll want to know: \bn{শোধরানো}}
\textbf{\bn{শোধরানো}} — \emph{shodhrāno} — \textquotedblleft{}to mend one's ways; to put right\textquotedblright{}.

\begin{quote}
\bn{শাস্তির} \bn{পরে} \bn{লোভী} \bn{রাজা} \bn{শোধরালেন।}
\end{quote}

\begin{quote}
\emph{shāstir pôre lobhī rājā shodhrālen.}
\end{quote}

\begin{quote}
After his punishment the greedy king mended his ways.
\end{quote}

Past (he/she …): \bn{সেই} \bn{ছেলে} \bn{শোধরাল}, \emph{sei chhele shodhrālô}.

\begin{cousinweb}
\textbf{\bn{শোধরানো}} (\emph{shodhrāno}) is usually explained as built on \bn{শোধ}, 'a clearing, a setting right' — Sanskrit \emph{śodha}, 'purification, correction' (Monier-Williams) — as a verb: to put oneself right.
\end{cousinweb}

\begin{grammarlens}[title={What you've built}]
One more word for telling what happened.
\end{grammarlens}

\subsection*{Your turn}
\begin{itemize}
\raggedright
  \item \emph{Say it:} \emph{shodhrāno}
  \item \emph{Say it:} \emph{shodhrāno}, once more
  \item \emph{Say it:} \emph{shāstir pôre lobhī rājā shodhrālen.} — After his punishment the greedy king mended his ways.
  \item \emph{Recall:} say the Bengali for an outcome, then the Bengali for a sin, then say \emph{shodhrāno} again
\end{itemize}

\begin{tcolorbox}[breakable,colback=teal!4,colframe=teal!35!black,title={Before you move on}]
What does \bn{শোধরানো} mean? (\textquotedblleft{}To mend one's ways; to put right\textquotedblright{}.) Say it once more.
\end{tcolorbox}
\section[\texorpdfstring{\bn{বাধ্য} \bn{হওয়া} (bādhyô hôoyā)}{বাধ্য হওয়া (bādhyô hôoyā)}]{\bn{বাধ্য} \bn{হওয়া} (bādhyô hôoyā) — to be forced; to have no choice}

\label{lesson:BN-C325-badhyo-hooya}



\begin{quote}
Before the new one: say the Bengali for a sin, then the Bengali for to mend one's ways.
\end{quote}

\subsection*{You'll want to know: \bn{বাধ্য} \bn{হওয়া}}
\textbf{\bn{বাধ্য} \bn{হওয়া}} — \emph{bādhyô hôoyā} — \textquotedblleft{}to be forced; to have no choice\textquotedblright{}.

\begin{quote}
\bn{ঝড়ের} \bn{জন্য} \bn{জেলেরা} \bn{ফিরে} \bn{আসতে} \bn{বাধ্য} \bn{হল।}
\end{quote}

\begin{quote}
\emph{jhôṛer jônyô jelerā phire āste bādhyô hôlô.}
\end{quote}

\begin{quote}
Because of the storm the fishermen were forced to come back.
\end{quote}

Past (he/she …): \bn{তিনি} \bn{বাড়ি} \bn{বেচতে} \bn{বাধ্য} \bn{হলেন}, \emph{tini bāṛi bechte bādhyô hôlen}.

\begin{cousinweb}
\textbf{\bn{বাধ্য} \bn{হওয়া}} (\emph{bādhyô hôoyā}) puts Sanskrit \emph{bādhya}, 'to be pressed hard, held back' (Monier-Williams), before \textbf{\bn{হওয়া}}: to be pressed into something. It takes the -\bn{তে} form: \bn{যেতে} \bn{বাধ্য} \bn{হল}, had to go.
\end{cousinweb}

\begin{grammarlens}[title={What you've built}]
One more word for telling what happened.
\end{grammarlens}

\subsection*{Your turn}
\begin{itemize}
\raggedright
  \item \emph{Say it:} \emph{bādhyô hôoyā}
  \item \emph{Say it:} \emph{bādhyô hôoyā}, once more
  \item \emph{Say it:} \emph{jhôṛer jônyô jelerā phire āste bādhyô hôlô.} — Because of the storm the fishermen were forced to come back.
  \item \emph{Recall:} say the Bengali for a sin, then the Bengali for to mend one's ways, then say \emph{bādhyô hôoyā} again
\end{itemize}

\begin{tcolorbox}[breakable,colback=teal!4,colframe=teal!35!black,title={Before you move on}]
What does \bn{বাধ্য} \bn{হওয়া} mean? (\textquotedblleft{}To be forced; to have no choice\textquotedblright{}.) Say it once more.
\end{tcolorbox}
